\begin{lemma}\label{lem:octic-quotient-cusps}
Let $\pi:X\longrightarrow Y$ be a finite morphism of smooth projective
curves over $\mathbb Q$, taking cusps to cusps, and let $c_0$ be a
rational cusp on $Y$. Suppose $\operatorname{Jac}(Y)(\mathbb Q)$ is
generated by rational cuspidal divisor classes. Let $p$ be a common
prime of good reduction and let $H_p\subseteq\operatorname{Jac}(Y)(\mathbb F_p)$
be the image of these global cuspidal classes. Write $q_p$ for the
quotient by $H_p$, and put
\[
 S_k=\{q_p([B-kc_0]): B\text{ effective cuspidal on }Y_{\mathbb F_p},
 \ \deg B=k\}.
\]
If an effective rational divisor of degree $d$ on $X$ reduces to
$D'+B'$, with $B'$ cuspidal and $\deg D'=d'$, then
\[
 -q_p([\pi_*D'-d'c_0])\in S_{d-d'}.
\]
\end{lemma}
\begin{proof}
The global class $[\pi_*D-dc_0]$ is rational, hence cuspidal by the
hypothesis. Reduction and proper pushforward commute. Both operations
preserve the degree of a divisor; in particular a closed point $P$
with image $Q$ pushes forward to $[\kappa(P):\kappa(Q)]Q$.
The image of $[\pi_*D-dc_0]$ in the quotient is therefore zero, and
$\pi_*B'$ is an effective cuspidal divisor of degree $d-d'$.
This gives the asserted membership. Allowing every finite-field
cuspidal complement on $Y$ gives a necessary condition, even when
some of them do not lift to global divisors.
\end{proof}
